\chapter{The runbook commands, flag by flag}
\label{ch:runbook}

\code{RUNBOOK.md} lists the commands that work. This chapter says what each
one does, so that when one fails you know which layer failed.

\section{One-time setup}

\begin{lstlisting}[style=shell]
winget install --id Python.Python.3.12 -e
py -0
\end{lstlisting}
Installs CPython 3.12 from the Windows package manager and lists the
installed interpreters (\code{-V:3.12} must appear). 3.10 fails twice over:
Newton's type annotations use \code{X | None} at run time, and the only
\code{imgui-bundle} wheel for 3.10 on Windows is too old, so pip tries to
compile it from source and fails.

\begin{lstlisting}[style=shell]
git clone -b develop https://github.com/isaac-sim/IsaacLab.git
py -3.12 -m venv .venv
Set-ExecutionPolicy -Scope Process -ExecutionPolicy Bypass
.\.venv\Scripts\Activate.ps1
python -m pip install --upgrade pip
.\isaaclab.bat -i
.\isaaclab.bat -i isaacsim
\end{lstlisting}
\begin{itemize}[leftmargin=1.6em]
  \item \code{-b develop}: the branch with the Newton backend and the
        \code{isaaclab train} CLI. \code{main} does not have them.
  \item \code{py -3.12 -m venv .venv}: a virtual environment inside the Isaac
        Lab folder; \code{isaaclab.bat} looks for it there.
  \item \code{Set-ExecutionPolicy ... Bypass}: lets PowerShell run the
        activation script in this window only.
  \item \texttt{.\textbackslash isaaclab.bat -i}: installs Isaac Lab's source packages
        (\code{isaaclab}, \code{isaaclab_newton}, \code{isaaclab_rl},
        \code{isaaclab_tasks}, \code{isaaclab_contrib}, \ldots) in editable
        mode plus Newton, Warp, torch and \code{rsl_rl}. This is the
        ``kit-less'' install: no Isaac Sim application.
  \item \code{-i isaacsim}: adds the Isaac Sim pip package, needed
        \emph{only} by the MJCF converter (Chapter~\ref{ch:asset}). Training
        never starts it.
\end{itemize}

\begin{lstlisting}[style=shell]
.\isaaclab.bat train --rl_library rsl_rl --task Isaac-Cartpole-Direct --num_envs 16 --max_iterations 10 physics=newton_mjwarp --viz newton
\end{lstlisting}
A smoke test with a built-in task before any of our code is involved.
\srcpath{physics=newton_mjwarp} is a Hydra typed selector that replaces the
task's physics section with the Newton preset (Cartpole defaults to PhysX);
our task does not need it because \code{__post_init__} builds its own
\code{NewtonCfg}. The pyglet \code{AssertionError} printed after
``Training time'' is the viewer window being torn down while a redraw is
pending; it is harmless.

\begin{lstlisting}[style=shell]
.\isaaclab.bat -p -m pip install -e C:\Users\hanse\Documents\Luna_HiFi
.\isaaclab.bat -p -c "import luna_hifi_tasks, gymnasium as gym; print([k for k in gym.registry if 'Luna-' in k])"
\end{lstlisting}
\code{-p} runs the venv's Python with the arguments that follow. The first
line installs our package in editable mode (Chapter~\ref{ch:pyproject}); the
second imports it and prints every registry key containing \code{Luna-}. The
expected output is \code{['Luna-Tricycle-Direct']}. An empty list means the
import ran but registered nothing, which is the missing-\code{__init__.py}
symptom; an \code{ImportError} means the editable install did not take.

\begin{lstlisting}[style=shell]
mkdir C:\Users\hanse\Documents\Luna_HiFi\assets
.\isaaclab.bat -p -c "from luna_hifi_tasks.tricycle.tricycle import write_mjcf; print(write_mjcf())"
.\isaaclab.bat -p scripts\tools\convert_mjcf.py C:\Users\hanse\AppData\Local\Temp\luna_hifi_assets\tricycle.xml C:\Users\hanse\Documents\Luna_HiFi\assets\tricycle.usd
\end{lstlisting}
The conversion. Line two calls the function from Chapter~\ref{ch:tricycle_py}
and prints the temp path it wrote. Line three runs Isaac Lab's converter
script on that file: positional arguments are the MJCF input and the USD
output. The converter ignores the exact output filename and writes
\texttt{assets\textbackslash tricycle\textbackslash tricycle.usda} plus the payload files
(Chapter~\ref{ch:usd}); the env cfg points at that real location. Run these
two lines again whenever \code{tricycle.py} changes.

\section{Every session}

\begin{lstlisting}[style=shell]
cd C:\Users\hanse\Documents\IsaacLab
Set-ExecutionPolicy -Scope Process -ExecutionPolicy Bypass
.\.venv\Scripts\Activate.ps1
\end{lstlisting}
\code{isaaclab.bat} must be run from the Isaac Lab folder (hence the
\texttt{.\textbackslash } prefix and the \code{cd}), and the venv must be active in each new
window; ``\code{No module named 'isaaclab_rl'}'' means it is not.

\begin{lstlisting}[style=shell]
.\isaaclab.bat train --rl_library rsl_rl --task Luna-Tricycle-Direct --num_envs 4
\end{lstlisting}
Chapter~\ref{ch:pipeline} end to end. Headless is the default; there is no
\code{--headless} flag on develop. Add \code{--viz newton} to watch, and
\code{--max_iterations N} to override the 300 in the agent cfg. Checkpoints
and TensorBoard logs go to \texttt{logs\textbackslash rsl\_rl\textbackslash tricycle\textbackslash <timestamp>\textbackslash }.

\begin{lstlisting}[style=shell]
.\isaaclab.bat play --rl_library rsl_rl --task Luna-Tricycle-Direct --num_envs 1 --checkpoint latest --viz newton
\end{lstlisting}
Section~\ref{sec:play}. \code{--checkpoint latest} picks the newest run's
last saved model; a path to a \code{model_N.pt} also works. \code{--num_envs
1} keeps the soil grid small; \code{--viz newton} opens the window in which
particles are drawn because of our \code{play_mode()}.

\section{Troubleshooting table, mapped to the code}

\begin{center}
\small
\begin{tabular}{@{}p{5.2cm}p{9.6cm}@{}}
\toprule
Symptom & Where in this manual \\
\midrule
\code{TypeError: Union[arg, ...]} on newton import & Python 3.10; Section~1 above. \\
\srcpath{NameNotFound} on the task id & Package not installed, or \code{pyproject.toml} edited without reinstalling: Chapter~\ref{ch:pyproject}, Section~\ref{sec:step_cli}. \\
gym registry empty after import & Missing \code{__init__.py}; add it, no reinstall: Chapters~\ref{ch:pyproject}, \ref{ch:agents}. \\
\code{A prim already exists at path} & Assets constructed in \code{_setup_scene}: Chapter~\ref{ch:env}. \\
capacity exceeded for upper/lower/leaf & Raise that cap, keep the ordering: cfg lines 257--260, Section~\ref{sec:mpm}. \\
\code{Failed to create volume} + CUDA 700 & Caps too high for memory; one failure, then the context is poisoned: Section~\ref{sec:mpm}. \\
\code{Mean episode length: 1.00} & A termination fires on step 1: check the flip sign and drift threshold, Chapter~\ref{ch:env}. \\
reward explodes to $10^8$ & Physics diverged; the invalid-state termination and reward clamp exist for this: Chapter~\ref{ch:env}. \\
soil invisible & \code{show_particles} defaults to \code{False}: cfg line 295. \\
car looks like the old box & Stale USD; reconvert: Chapter~\ref{ch:asset}. \\
\bottomrule
\end{tabular}
\end{center}

After any CUDA error 700, check \code{nvidia-smi} and kill leftover
\code{python.exe} processes before retrying: the CUDA context is unusable and
the VRAM may still be held.
